The next theorem concerns compactness of imbeddings between Sobolev spaces.

\begin{Theorem}
\label{thm:compact-successive-sobolev-orders}
Let $\Omega$ be bounded and let $k$ be a non-negative integer. Assume that $\Omega$ has the $(k+1)$-extension property. Then\[H^{k+1}(\Omega)\overset{c}{\hookrightarrow}H^k(\Omega),\]i.e., $H^{k+1}(\Omega)$ is compactly imbedded in $H^k(\Omega)$.
\end{Theorem}

\begingroup
\postdisplaypenalty=10000
\renewcommand{\endProof}{\par\nobreak\hfill$\blacksquare$\par\vspace{0.1cm}}
\begin{Proof}
Let $E$ be an extension operator; we can always choose $E$ in such a way that it maps to functions supported in some compact set $S$. If $u_n$ is a bounded sequence in $H^{k+1}(\Omega)$, then $Eu_n$ is a bounded sequence in $H^{k+1}(\mathbb{R}^m)$ and the support of $Eu_n$ is contained in $S$. We shall show that $Eu_n$ has a subsequence which converges in $H^k(\mathbb{R}^m)$. Let $M_R$ denote the operator of multiplication by the characteristic function of the ball $|\bm{\xi}|<R$. For $v\in H^{k+1}(\mathbb{R}^m)$, let $v^R=\mathcal{F}^{-1}[M_R\mathcal{F}[v]]$. Then it is easy to show that there is a constant $C$ with
\begin{equation}\label{eq:sobolev-fourier-cutoff-high-frequency-bound}
\|u-v^R\|_{k,2}\leq\dfrac{C}{R}\|u\|_{k+1,2}.
\end{equation}
Moreover for any fixed $R$ and any positive integer $l$, we have $u^R\in C_b^l(\mathbb{R}^m)$ and there is a constant $C(R,l)$ such that
\begin{equation}\label{eq:sobolev-fourier-cutoff-smooth-norm-bound}
\|u^R\|_{l,\infty}\leq C(R,l)\|u\|_{k+1,2}.
\end{equation}
Hence if $Eu_n$ is bounded in $H^{k+1}(\mathbb{R}^m)$, and $D$ is any bounded open set containing $S$, then $(Eu_n)^R$ has a subsequence which converges in $C_b^k(\overline D)$ (Arzela-Ascoli theorem) and hence in $H^k(D)$. By choosing a sequence $R_n\to\infty$ and applying a diagonal argument, we find a subsequence $n_j$ such that $(Eu_{n_j})^{R_n}$ converges in $H^k(D)$ for every $n$. In conjunction with \eqref{eq:sobolev-fourier-cutoff-high-frequency-bound}, this implies that $Eu_{n_j}$ converges in $H^k(D)$ and hence in $H^k(\mathbb{R}^m)$, since the support of each $Eu_{n_j}$ is contained in $D$.
\end{Proof}
\endgroup

As a consequence of Theorem~\ref{thm:compact-successive-sobolev-orders}, we shall establish Ehrling\textquotesingle s lemma, which allows us to introduce an equivalent norm on $H^k(\Omega)$. We first give an abstract version.

\begin{Theorem}{Ehrling\textquotesingle s lemma}
\label{thm:ehrling-compact-intermediate-norm}
Let $X,Y$ and $Z$ be Banach spaces. Assume that $X$ is compactly imbedded in $Y$ and $Y$ is continuously imbedded in $Z$\[X\overset{c}{\hookrightarrow}Y\hookrightarrow Z.\]Then for every $\varepsilon>0$ there exists a constant $c(\varepsilon)$ such that
\begin{equation}\label{eq:ehrling-compact-intermediate-norm-bound}
\|x\|_Y\leq\varepsilon\|x\|_X+c(\varepsilon)\|x\|_Z
\end{equation}
for every $x\in X$.
\end{Theorem}

\begingroup
\postdisplaypenalty=10000
\renewcommand{\endProof}{\par\nobreak\hfill$\blacksquare$\par\vspace{0.1cm}}
\begin{Proof}
Assume the claim fails for some $\varepsilon_0>0$. Then there is a sequence $x_n$ in $X$ such that $\|x_n\|_X=1$ and
\begin{equation}\label{eq:ehrling-normalized-contradiction-sequence}
\|x_n\|_Y>\varepsilon_0+n\|x_n\|_Z.
\end{equation}
Since the imbedding from $X$ to $Y$ is continuous, $x_n$ is bounded in $Y$, and \eqref{eq:ehrling-normalized-contradiction-sequence} implies that $x_n$ must converge to $0$ in $Z$. After passing to a subsequence, we may assume that $x_n$ converges in $Y$; the limit must then be $0$. But this contradicts \eqref{eq:ehrling-normalized-contradiction-sequence}.
\end{Proof}
\endgroup
